Une bobine rectangulaire de $N = \v{N}$ spires mesure \qty{\v{a}}{cm} sur \qty{\v{b}}{cm}. Elle est parcourue par un courant de \qty{\v{I}}{A} et placée dans un champ magnétique uniforme de \qty{\v{B}}{T}. La normale au plan de la bobine fait un angle de \qty{\v{th}}{\degree} avec le champ.

\begin{enumerate}[label=\alph*)]
    \item Déterminez le moment magnétique de la bobine.
    \item Déterminez le moment de force exercé sur la bobine.
    \item Déterminez le moment de force maximal que ce champ peut exercer sur la bobine.
    \item Quel travail un agent extérieur doit-il fournir pour faire tourner la bobine de la position où sa normale est parallèle au champ jusqu'à la position décrite dans l'énoncé ?
\end{enumerate}
